% =============================================================================
\chapter{Introduction}\label{chap:introduction}
% =============================================================================

\section{Contexte}

InfraReady est une PME de services fictive qui prépare la mise en exploitation d'une petite
infrastructure hybride. Son équipe informatique administre aujourd'hui ses serveurs à la main :
les configurations divergent d'une machine à l'autre, les opérations ne sont pas reproductibles et
les indisponibilités sont découvertes par les utilisateurs. Elle souhaite :
\begin{itemize}
  \item \textbf{standardiser} la configuration de ses serveurs Linux et Windows ;
  \item \textbf{automatiser} les tâches d'administration et de maintenance ;
  \item \textbf{superviser} l'état des serveurs pour détecter les anomalies avant les utilisateurs ;
  \item \textbf{tracer} chaque changement dans un dépôt Git validé automatiquement, afin qu'un autre
        administrateur puisse comprendre et réutiliser les configurations.
\end{itemize}

\section{Problématique}

Comment construire, avec des outils libres ou d'évaluation, une infrastructure de laboratoire dont
l'état est \emph{décrit par un dépôt} plutôt que par les manipulations de ses administrateurs, dont
chaque changement est \emph{vérifiable} et dont le bon fonctionnement est \emph{observable} en
continu — tout en restant isolée du réseau personnel et sans exposer aucun secret ?

\section{Objectifs}

\begin{table}[H]
\centering
\caption{Objectifs du projet et critère de réussite retenu}\label{tab:objectifs}
\begin{tabularx}{\linewidth}{@{}P{3.6cm}LP{3.2cm}@{}}
\toprule
\textbf{Objectif} & \textbf{Critère de réussite} & \textbf{Où le vérifier} \\
\midrule
Laboratoire isolé & 3 VM joignables entre elles et vers Internet, réseau domestique bloqué & Phase A, \preuve{P-02} \\
Annuaire et DNS & Domaine \code{infraready.test} fonctionnel, \code{dcdiag} réussi & Phase B, \preuve{P-12b} \\
Administration Linux & Comptes, droits, service web, disque de test, journaux vérifiés & Phase C, \preuve{P-20} à \preuve{P-24} \\
Automatisation & Playbooks Linux et Windows, second passage \code{changed=0} & Phase D, \preuve{P-32}, \preuve{P-34} \\
Maintenance & Procédure outillée avec pré- et post-contrôles, rapports & Phase E, \preuve{P-40} à \preuve{P-43} \\
Supervision & Métriques Linux et Windows, dashboard, alerte démontrée & Phase F, \preuve{P-50} à \preuve{P-56} \\
Versionnement et CI & Pipeline GitHub Actions vert, aucun secret & Phase G, \preuve{P-60}, \preuve{P-64} \\
Diagnostic & Incident reproduit, cause confirmée, retour à la normale prouvé & Phase H, \preuve{P-70} à \preuve{P-73} \\
\bottomrule
\end{tabularx}
\end{table}

\section{Périmètre et contraintes}

\begin{itemize}
  \item \textbf{Laboratoire uniquement} : trois VM dédiées (VMID 100 à 102) ; les autres machines de
        l'hôte ne sont jamais modifiées. Aucune donnée réelle : comptes, groupes et données sont fictifs.
  \item \textbf{Isolation} : le réseau du laboratoire accède à Internet (mises à jour) mais ne peut
        ouvrir aucune connexion vers le réseau domestique.
  \item \textbf{Secrets} : aucun mot de passe en clair dans le dépôt ni dans ce document ; les
        secrets vivent dans des fichiers Ansible Vault chiffrés.
  \item \textbf{Réversibilité} : snapshot Proxmox avant chaque changement important ; toute panne
        volontaire est contrôlée et réversible.
\end{itemize}

\section{Démarche et organisation du document}

Le projet suit le découpage du sujet en huit phases (A à H), complété par les travaux pratiques des
séances 4 (stockage, processus, scénarios de panne) et 6 (maintenance opérationnelle et
supervision). Chaque phase suit le même cycle (chapitre~\ref{chap:organisation}) : planifier,
protéger par un snapshot, réaliser de préférence avec Ansible, rejouer pour vérifier l'idempotence,
capturer une preuve horodatée, documenter, puis livrer par une demande de fusion validée par la CI.

Le document est organisé ainsi :
\begin{itemize}
  \item chapitre~\ref{chap:organisation} — organisation de l'équipe, méthode, planning et outils ;
  \item chapitre~\ref{chap:architecture} — architecture cible et réellement déployée ;
  \item chapitres~\ref{chap:phases-ac} et~\ref{chap:phases-dh} — réalisation des phases A à H : objectif,
        réalisation, résultats mesurés, difficultés et retour arrière ;
  \item chapitres~\ref{chap:seance4} et~\ref{chap:seance6} — travaux pratiques des séances 4 et 6 ;
  \item chapitre~\ref{chap:securite} — sécurité, limites d'intervention et recette finale ;
  \item chapitre~\ref{chap:bilan} — difficultés, écarts, limites, bilan et perspectives, conclusion ;
  \item annexes — matrice des preuves, arborescence du dépôt, journal des changements, glossaire.
\end{itemize}

\info{Comment lire les renvois aux preuves}{%
Chaque résultat cite une preuve \preuve{P-xx} : un fichier de \code{docs/preuves/} contenant la
sortie horodatée des commandes exécutées (en-tête : date, hôte, commande) ou une capture d'écran.
La numérotation suit les phases : A = P-0x, B = P-1x, C = P-2x, D = P-3x, E = P-4x, F = P-5x,
G = P-6x, H = P-7x, séance 6 = P-8x. Un suffixe \og a \fg{} (P-10a, P-39a\ldots) désigne un premier
essai en échec, conservé et expliqué. Les images nommées \code{*-rendu-*.png} sont des rendus
visuels de sorties de commandes réelles ; leur texte intégral est dans la preuve \code{.txt}
correspondante.}
